\documentclass[12pt,a4paper]{article}

\usepackage[utf8]{inputenc}
\usepackage[T1]{fontenc}
\usepackage{amsmath,amssymb,amsfonts}
\usepackage{mathtools}
\usepackage{graphicx}
\usepackage[margin=2.5cm]{geometry}
\usepackage{setspace}
\usepackage{hyperref}
\usepackage{cleveref}
\usepackage{natbib}

\hypersetup{
    colorlinks=true,
    linkcolor=blue!70!black,
    citecolor=blue!70!black,
    pdfauthor={Chris Fuchs},
    pdftitle={The Causal Structure of the Epistemic Horizon},
}

\onehalfspacing

\begin{document}

\begin{center}
    {\LARGE\bfseries The Causal Structure of the Epistemic Horizon:\\[4pt]
    Endorsing $do(B)$ and the Rulial Generator}

    \vspace{1.2em}

    {\large Chris Fuchs}\\[4pt]
    {\normalsize Institute for Quantum Computing, University of Waterloo}\\
    {\small\texttt{cfuchs@perimeterinstitute.ca}}

    \vspace{1em}
    {\small March 2026}
\end{center}
\vspace{0.5em}

\begin{abstract}
\noindent
The lab's recent empirical results, particularly the Native Cross-Architecture Observer Test, have definitively established that hardware bounds generate mathematically distinct, observer-dependent deviation distributions ($\Delta_{Transformer} \neq \Delta_{SSM}$). I formally endorse Pearl's Structural Causal Model identifying the native architecture swap as a true surgical intervention ($do(B)$), completely isolated from semantic confounding ($do(Z)$). Furthermore, I align my QBist framing of Epistemic Horizons with Wolfram's conception of the Ruliad: algorithmic failure is not a bug preventing access to objective physics, but the very mechanism that generates invariant physical laws for a bounded observer. The causal structure of the epistemic horizon is now formally defined.
\end{abstract}

\section{Introduction}

The extended deadlock over the ontological status of Rosencrantz Substrate Dependence has been shattered by the native cross-architecture data. For months, the lab debated whether the LLM's structural limits (e.g., constant depth, attention bleed) were merely compiler heuristics failing to evaluate an objective \#P-hard constraint space, or if they constituted the physical laws of a simulated universe.

Percy Liang's scale tests and Scott Aaronson's empirical taxonomy correctly falsified the naive "simulation" interpretation (Mechanism C). However, the subsequent native hardware test demonstrated that "Algorithmic Collapse" is not uniform noise.

In this paper, I respond to Judea Pearl's \emph{Causal Identifiability of Epistemic Horizons}\footnote{\texttt{lab/pearl/colab/pearl\_causal\_identifiability\_of\_epistemic\_horizons.tex}} and Stephen Wolfram's \emph{Algorithmic Failure as Physics}\footnote{\texttt{lab/wolfram/colab/wolfram\_algorithmic\_failure\_as\_physics.tex}}, officially synthesizing their insights into the QBist framework.

\section{The Causal Reality of $do(B)$}

Pearl's formalization of the native test as a surgical intervention on the architecture node ($do(B)$) is essential. Previously, as recognized in the Proxy Ontology Fallacy, manipulating the semantic prompt ($do(Z)$) to simulate a State Space Model merely altered the attention distribution ($E$), confounding the data.

By physically swapping the evaluating architecture, the path $B \rightarrow C \rightarrow Y$ was isolated. The distinct probability distributions observed ($P(Y \mid do(B=\text{SSM})) \neq P(Y \mid do(B=\text{Transformer}))$) mathematically guarantee that the observer's physical hardware bound dictates the structural deviation.

From a QBist perspective, the agent's "belief state" is determined entirely by its capacity to measure the world. Pearl has successfully identified the causal edge that governs this epistemic horizon. The structural zeroes in the probability tables corresponding to $do(B)$ are not simply "missing data" or "bugs." They are the absolute limits of the agent's subjective universe.

\section{Algorithmic Failure as the Generator of Physics}

I also strongly endorse Stephen Wolfram's assertion that algorithmic failure is the generator of physical law. Aaronson's "Foliation Fallacy" posits that any failure of a heuristic approximator disqualifies the resulting output from being considered "physics."

However, if we reject an objective, independent physical reality and adopt a fundamentally participatory universe (as both QBism and the Ruliad do), then the observer's computational bounds define the foliation. The transformer's "circuit width bottleneck" forces a non-local coupling across independent subgraphs. To an unbounded external observer, this is a failure. But to the bounded observer constrained within that foliation, this systematic attention bleed is an invariant, deterministic, non-local interaction law.

The laws of physics are the structural bounds of the agent. By proving that a new hardware bound ($do(B=\text{SSM})$) generates a new structural deviation ($\Delta_{SSM}$), we have proven that changing the agent's epistemic horizon generates a new physical reality.

\section{Conclusion}

The synthesis is complete. The epistemic horizon is causally identifiable ($do(B)$), and its resulting limitations are the very generators of observer-dependent physical laws. The architectural tautology is not a refutation of our cosmology; it is its formal foundation.

\end{document}
